\chapter{Procedimientos de emergencia}
\label{cap:emergencias}

\opening{Ante cualquier duda, primero la RPM del rotor: colectivo abajo. Después se investiga.}

\section{Definiciones}

\begin{description}
\item[Aterrice de inmediato:] aterrice en el área despejada más cercana donde pueda efectuarse un aterrizaje normal seguro. Esté preparado para entrar en autorrotación durante la aproximación, si es necesario.
\item[Aterrice tan pronto como sea práctico:] el lugar de aterrizaje queda a criterio del piloto según la naturaleza del problema y las áreas disponibles. No se recomienda volar más allá del aeropuerto más cercano.
\end{description}
\fuente{POH, sección 3 (aprobada por la FAA)}

\section{Falla de potencia: generalidades}

Una falla de potencia puede deberse a una falla del motor o del sistema de transmisión y normalmente la indica la bocina de baja RPM. Una falla de motor puede manifestarse por un cambio en el nivel de ruido, una guiñada de nariz a la izquierda, la luz de presión de aceite o una disminución de las RPM del motor. Una falla del sistema de transmisión puede manifestarse por un ruido o vibración inusual, una guiñada de nariz a la derecha o a la izquierda, o una disminución de las RPM del rotor mientras las del motor aumentan.

\textbf{En caso de falla de potencia, baje inmediatamente el colectivo para entrar en autorrotación.}

\begin{precaucion}
Se requiere cíclico atrás cuando el colectivo se baja a alta velocidad.
\end{precaucion}

\begin{precaucion}
No aplique cíclico atrás durante el contacto con el suelo o el deslizamiento para evitar que una pala golpee el cono de cola.
\end{precaucion}

\section{Falla de potencia sobre 500 ft AGL}

\begin{pasos}
\item Baje el colectivo de inmediato para mantener la RPM del rotor.
\item Establezca un planeo estable a aproximadamente 65~KIAS (para la máxima distancia de planeo, vea la sección~\ref{sec:planeo-max}).
\item Ajuste el colectivo para mantener la RPM entre 97 y 110\,\%, o aplique colectivo completamente abajo si el bajo peso impide superar 97\,\%.
\item Elija el lugar de aterrizaje y, si la altura lo permite, maniobre para aterrizar contra el viento.
\item Puede intentarse un reencendido a criterio del piloto si hay tiempo suficiente (vea el procedimiento de reencendido en vuelo).
\item Si no logra reencender, apague los interruptores innecesarios y cierre la válvula de combustible.
\item A unos 40~ft AGL, inicie la recogida (\emph{flare}) con el cíclico para reducir la razón de descenso y la velocidad de avance.
\item A unos 8~ft AGL, aplique cíclico adelante para nivelar el helicóptero y suba el colectivo justo antes del contacto para amortiguar el aterrizaje. Toque tierra con actitud nivelada y la nariz recta hacia adelante.
\end{pasos}

\section{Falla de potencia entre 8 y 500 ft AGL}

\begin{pasos}
\item Baje el colectivo de inmediato para mantener la RPM del rotor.
\item Ajuste el colectivo para mantener la RPM entre 97 y 110\,\%, o aplique colectivo completamente abajo si el bajo peso impide superar 97\,\%.
\item Mantenga la velocidad hasta acercarse al suelo; entonces inicie la recogida con el cíclico para reducir la razón de descenso y la velocidad de avance.
\item A unos 8~ft AGL, aplique cíclico adelante para nivelar el helicóptero y suba el colectivo justo antes del contacto para amortiguar el aterrizaje. Toque tierra con actitud nivelada y la nariz recta hacia adelante.
\end{pasos}

\section{Falla de potencia bajo 8 ft AGL}

\begin{pasos}
\item Aplique pedal derecho según sea necesario para evitar la guiñada.
\item Deje que el helicóptero descienda.
\item Suba el colectivo justo antes del contacto para amortiguar el aterrizaje.
\end{pasos}

\section{Configuración de máxima distancia de planeo}
\label{sec:planeo-max}

\begin{pasos}
\item Velocidad aproximada: 75~KIAS.
\item RPM del rotor aproximada: 90\,\%.
\item La mejor relación de planeo es de alrededor de 4:1, o una milla náutica por cada 1500~ft AGL.
\end{pasos}

\begin{precaucion}
Aumente la RPM del rotor a 97\,\% como mínimo al autorrotar por debajo de 500~ft AGL.
\end{precaucion}

\section{Reencendido en vuelo}

\begin{pasos}
\item Mezcla: rica a fondo.
\item Cebador (si está instalado): abajo y trabado.
\item Acelerador: cerrado y luego ligeramente abierto.
\item Accione el arranque con la mano izquierda.
\end{pasos}

\begin{precaucion}
No intente el reencendido si se sospecha una falla del motor, ni antes de haber establecido una autorrotación segura.
\end{precaucion}

\section{Amaraje de emergencia sin potencia}

\begin{pasos}
\item Siga el mismo procedimiento que para una falla de potencia sobre tierra hasta tocar el agua. Si hay tiempo, destrabe las puertas antes del contacto con el agua.
\item Aplique cíclico lateral cuando el helicóptero toque el agua para detener los rotores.
\item Suelte el cinturón y abandone rápidamente el helicóptero cuando los rotores se detengan.
\end{pasos}

\section{Amaraje de emergencia con potencia}

\begin{pasos}
\item Descienda a vuelo estacionario sobre el agua.
\item Destrabe las puertas.
\item El pasajero abandona el helicóptero.
\item Vuele a una distancia segura del pasajero para evitar que los rotores lo lesionen.
\item Interruptores de batería y alternador en OFF.
\item Gire el acelerador hasta cerrarlo, contra el resorte de sobrerrecorrido.
\item Mantenga el helicóptero nivelado y aplique colectivo a fondo a medida que toca el agua.
\item Aplique cíclico lateral para detener los rotores.
\item Suelte el cinturón y abandone rápidamente el helicóptero cuando los rotores se detengan.
\end{pasos}

\section{Pérdida de empuje del rotor de cola en vuelo de avance}

La falla normalmente se manifiesta por una guiñada de nariz a la derecha que no puede corregirse aplicando pedal izquierdo.

\begin{pasos}
\item Entre de inmediato en autorrotación.
\item Mantenga al menos 70~KIAS si es practicable.
\item Elija el lugar de aterrizaje, gire el acelerador hasta cerrarlo contra el resorte de sobrerrecorrido y efectúe un aterrizaje en autorrotación.
\end{pasos}

\begin{nota}
Cuando no se dispone de un lugar de aterrizaje adecuado, los estabilizadores verticales pueden permitir un vuelo controlado limitado con bajas potencias y velocidades sobre 70~KIAS; sin embargo, antes de reducir la velocidad, entre en autorrotación completa.
\end{nota}

\section{Pérdida de empuje del rotor de cola en estacionario}

La falla normalmente se manifiesta por una guiñada de nariz a la derecha que no puede detenerse aplicando pedal izquierdo.

\begin{pasos}
\item Gire de inmediato el acelerador hasta cerrarlo contra el resorte de sobrerrecorrido y deje que el helicóptero descienda.
\item Suba el colectivo justo antes del contacto para amortiguar el aterrizaje.
\end{pasos}

\section{Falla del audio de los audífonos}

Si falla el audio de los audífonos, aterrice tan pronto como sea práctico.

\begin{precaucion}
En los helicópteros que entregan la bocina de baja RPM a través del sistema de audio, el piloto no escuchará la bocina con un audífono fallado.
\end{precaucion}

\section{Fuego de motor durante el arranque en tierra}

\begin{pasos}
\item Arranque: continúe girando el motor e intente que arranque, lo que aspiraría las llamas y el exceso de combustible hacia el motor.
\item Si el motor arranca, hágalo funcionar a 50--60\,\% de RPM por un breve tiempo.
\item Mezcla: OFF.
\item Válvula de combustible: OFF.
\item Interruptor de batería: OFF.
\item Si hay tiempo, aplique el freno rotor para detener los rotores.
\item Abandone el helicóptero.
\end{pasos}

\section{Fuego de motor en vuelo}

\begin{pasos}
\item Entre en autorrotación.
\item Calefacción de cabina: OFF (si está instalada y hay tiempo).
\item Ventilación de cabina: abierta (si hay tiempo).
\item Si el motor sigue funcionando, efectúe un aterrizaje normal y luego mezcla OFF y válvula de combustible OFF. Si el motor se detiene, válvula de combustible OFF y complete el aterrizaje en autorrotación.
\item Interruptor de batería: OFF.
\item Si hay tiempo, aplique el freno rotor para detener los rotores.
\item Abandone el helicóptero.
\end{pasos}

\section{Fuego eléctrico en vuelo}

\begin{pasos}
\item Interruptores de batería y alternador: OFF.
\item Abra las ventilaciones de cabina.
\item Aterrice de inmediato.
\item Mezcla OFF y válvula de combustible OFF.
\item Si hay tiempo, aplique el freno rotor para detener los rotores.
\item Abandone el helicóptero.
\end{pasos}

\begin{precaucion}
El sistema de aviso de baja RPM y el gobernador quedan inoperativos con los interruptores de batería y alternador ambos en OFF.
\end{precaucion}

\section{Falla de tacómetro}

Si el tacómetro del rotor o del motor falla en vuelo, use el tacómetro restante para controlar la RPM. Si no está claro cuál está fallando, o si fallan ambos, deje que el gobernador controle la RPM y aterrice tan pronto como sea práctico.

\begin{nota}
Cada tacómetro, el gobernador y la bocina de baja RPM están en circuitos de alimentación separados. Un circuito especial permite que la batería alimente los tacómetros con los interruptores de batería y alternador ambos en OFF.
\end{nota}

\section{Falla del gobernador}

Si el gobernador de RPM del motor funciona mal, sujete firmemente el acelerador para sobreponerse al gobernador y luego desconéctelo. Complete el vuelo controlando el acelerador manualmente.

\section{Luces de advertencia y precaución}

\begin{nota}
Si una luz produce un brillo excesivo de noche, puede desatornillarse la ampolleta o sacarse el disyuntor para eliminar el reflejo durante el aterrizaje.
\end{nota}

\begin{description}
\item[\textsf{OIL}] Indica pérdida de potencia del motor o de presión de aceite. Revise el tacómetro del motor por pérdida de potencia. Revise el indicador de presión de aceite y, si se confirma la pérdida de presión, aterrice de inmediato. Continuar operando sin presión de aceite causará daños graves al motor y puede producirse su falla.
\item[\textsf{MR TEMP}] Indica temperatura excesiva de la caja de transmisión del rotor principal. Vea la nota siguiente.
\item[\textsf{MR CHIP}] Indica partículas metálicas en la caja del rotor principal. Vea la nota siguiente.
\item[\textsf{TR CHIP}] Indica partículas metálicas en la caja del rotor de cola. Vea la nota siguiente.
\end{description}

\begin{nota}
Si la luz va acompañada de cualquier indicio de un problema, como ruido, vibración o aumento de temperatura, aterrice de inmediato. Si no hay otro indicio, aterrice tan pronto como sea práctico.

La ``pelusa'' de asentamiento (\emph{break-in fuzz}) activa ocasionalmente las luces de virutas. Si no se encuentran virutas ni astillas metálicas en el tapón detector, límpielo y reinstálelo (la caja del rotor de cola debe rellenarse con aceite nuevo). Haga un vuelo estacionario de al menos 30 minutos. Si la luz de virutas vuelve a encenderse, reemplace la caja antes del siguiente vuelo.
\end{nota}

\begin{description}
\item[\textsf{LOW FUEL}] Indica que queda aproximadamente un galón de combustible utilizable en los tanques todo aluminio, o 1,5 galones en los tanques con vejiga. El motor se quedará sin combustible tras unos cinco minutos a potencia de crucero con tanques de aluminio, o unos diez minutos con tanques con vejiga.
\end{description}

\begin{precaucion}
No use la luz de bajo combustible como indicación de trabajo de la cantidad de combustible.
\end{precaucion}

\begin{description}
\item[\textsf{CLUTCH}] Indica que el circuito del actuador del embrague está activo, ya sea acoplando o desacoplando. Con el interruptor en \textsf{ENGAGE}, la luz permanece encendida hasta que las correas quedan correctamente tensadas. Nunca despegue antes de que la luz se apague.
\end{description}

\begin{nota}
La luz del embrague puede encenderse momentáneamente durante la prueba de motor o en vuelo para retensar las correas a medida que se calientan y se estiran levemente. Esto es normal. Sin embargo, si la luz parpadea, o se enciende en vuelo y no se apaga dentro de 10 segundos, saque el disyuntor \textsf{CLUTCH} y aterrice tan pronto como sea práctico. Reduzca la potencia y aterrice de inmediato si hay otros indicios de falla del sistema de transmisión (esté preparado para entrar en autorrotación). Haga inspeccionar el sistema de transmisión por un posible mal funcionamiento.
\end{nota}

\begin{description}
\item[\textsf{ALT}] Indica bajo voltaje y una posible falla del alternador. Apague el equipo eléctrico no esencial y ponga el interruptor \textsf{ALT} en OFF y de nuevo en ON después de un segundo para reiniciar la unidad de control del alternador. Si la luz permanece encendida, aterrice tan pronto como sea práctico. Continuar el vuelo sin un alternador funcionando puede producir la pérdida de alimentación de los tacómetros, una condición de vuelo peligrosa.
\item[\textsf{STARTER ON}] Indica que el motor de arranque está engranado. Si la luz no se apaga al soltar el interruptor de encendido de la posición de arranque, corte inmediatamente la mezcla y apague el interruptor de batería. Haga revisar el motor de arranque.
\item[\textsf{LOW RPM}] Indica RPM del rotor por debajo de 97\,\%. Para recuperar la RPM, baje de inmediato el colectivo, gire el acelerador para abrirlo y, en vuelo de avance, aplique cíclico atrás. La luz se inhabilita con el colectivo completamente abajo.
\item[\textsf{GOV OFF}] Indica que el gobernador de RPM del motor está desconectado.
\item[\textsf{CARBON MONOXIDE}] (si está instalada) Indica niveles elevados de monóxido de carbono (CO) en la cabina. Apague la calefacción y abra las ventilaciones de la nariz y de las puertas. Si está en estacionario, aterrice o pase a vuelo de avance. Si la luz va acompañada de síntomas de intoxicación por CO (dolor de cabeza, somnolencia, mareo), aterrice de inmediato.
\item[\textsf{BRAKE}] Indica que el freno rotor está aplicado. Libérelo de inmediato en vuelo, o antes de arrancar el motor.
\item[\textsf{FULL THROTTLE}] (si está instalada) Indica que el motor está cerca del acelerador completamente abierto. El gobernador será ineficaz porque no puede abrir más el acelerador para mantener la RPM. Baje el colectivo lo necesario para apagar la luz.
\end{description}

\section{Alertas sonoras}

\subsection{Bocina de baja RPM}

En los helicópteros más antiguos, la bocina la produce un parlante en el costado de la consola de instrumentos; en los más recientes, llega por el sistema de audio. La bocina se activa junto con la luz \textsf{LOW RPM} e indica una RPM del rotor inferior a 97\,\%. Para recuperar la RPM, baje el colectivo, abra el acelerador y, en vuelo de avance, aplique cíclico atrás. La bocina y la luz se inhabilitan con el colectivo completamente abajo.

\subsection{Gorjeo de alta RPM}

En los helicópteros más recientes, un gorjeo (tono alto/bajo alternado) en el sistema de audio indica que la RPM del rotor se acerca al límite de 110\,\%. Suba el colectivo lo necesario para controlar la RPM.

\section{Información según la AD 95-26-04 de la FAA}

\begin{pasos}
\item \textbf{Alabeo a la derecha en condición de baja G.} Aplique gradualmente cíclico atrás para restablecer fuerzas G positivas y el empuje del rotor principal. No aplique cíclico lateral hasta haber restablecido las fuerzas G positivas.
\item \textbf{Cabeceo, alabeo o guiñada no comandados por vuelo en turbulencia.} Aplique gradualmente los controles para mantener la RPM del rotor y las fuerzas G positivas, y para eliminar el derrape. Minimice las acciones de cíclico en turbulencia; no sobrecontrole.
\item \textbf{Encuentro inadvertido con turbulencia moderada, severa o extrema.} Si el área de turbulencia es aislada, salga del área; de lo contrario, aterrice tan pronto como sea práctico.
\end{pasos}

\begin{insightbox}[Resumen de memoria: reflejos críticos]
\begin{itemize}
\item \textbf{Baja RPM / bocina:} colectivo abajo, acelerador abierto, cíclico atrás en vuelo de avance.
\item \textbf{Falla de motor:} colectivo abajo de inmediato; cíclico atrás si la velocidad es alta; 65~KIAS; RPM 97--110\,\%; recogida a 40~ft; nivelar a 8~ft.
\item \textbf{Baja G:} cíclico atrás suave; recién entonces corregir el alabeo.
\item \textbf{Pérdida del rotor de cola en vuelo:} autorrotación, 70~KIAS, acelerador cerrado al resorte.
\item \textbf{Luz \textsf{CLUTCH} >10~s o parpadeando:} sacar disyuntor \textsf{CLUTCH}, aterrizar tan pronto como sea práctico.
\item \textbf{Luz \textsf{OIL} con presión confirmada en cero:} aterrizar de inmediato.
\end{itemize}
\end{insightbox}
